% Volume VI: Machine Learning Mathematics
% Continuity note for Chapter 33.
% Previous dependencies: Chapters 5-7 provide least squares, projections, norms, and matrix calculus; Chapters 12 and 18 provide function-space and information geometry intuition; Chapters 13-15 provide Gaussian conditioning and Bayes; Chapters 31-32 provide ERM, generalization, linear models, GLMs, and classification.
% New durable concepts: regularized risk, complexity penalties, ridge and lasso intuition, MAP interpretation, kernel functions, feature maps, positive semidefinite Gram matrices, kernel trick, kernel ridge regression, RKHS, reproducing property, representer theorem, Gaussian processes as distributions over functions, GP regression posterior, kernel hyperparameters, noise variance, uncertainty, and computational caveats.
% Forward hooks: Chapter 34 uses priors and posterior conditioning in latent-variable models; Chapter 35 reuses kernel and manifold ideas; Chapter 41 uses Gaussian noise, KL terms, and diffusion connections; Chapter 48 uses GPs and kernels for latent neural trajectories; Chapter 51 reuses regularization and predictive uncertainty for model comparison.
% Notation commitments: Regularized objectives have empirical loss plus lambda Omega; feature maps are phi:\mathcal X\to\mathcal H; kernels are k(x,x')=<phi(x),phi(x')>_\mathcal H; Gram matrices are K_{ij}=k(x_i,x_j); RKHS norm is ||f||_\mathcal H; Gaussian processes are f\sim\mathcal{GP}(m,k); observation noise variance is sigma_n^2; training kernel matrix is K; test-training kernel vector is k_*; posterior mean and covariance use K+sigma_n^2 I_n.
% Figure plan: no missing images are included. Proposed ImageGen prompts are recorded in imagegen-figures/ch33-figure-metadata.md.
\section{Chapter 33: Regularization, Kernels, and Gaussian Processes}
\label{ch:regularization-kernels-gaussian-processes}

Chapter 32 showed how linear scores support regression, classification, GLMs, and margins. Those methods already reveal a central tension: flexible models fit more patterns, but finite data cannot distinguish every pattern from noise. Chapter 33 studies three linked responses to that tension. Regularization encodes preferences among functions. Kernels let linear methods operate in rich feature spaces without explicitly constructing all features. Gaussian processes put probability distributions directly on functions and use Gaussian conditioning to produce predictions with uncertainty.

The central problem is to control function complexity while keeping enough flexibility to model nonlinear structure. The mathematical bridge is the norm of a function in a feature or kernel space. In a deterministic method, the norm appears as a regularizer. In a Bayesian method, the same structure appears as a prior covariance over functions.

\paragraph{Chapter problem.}
How can prior knowledge about smoothness, size, similarity, and uncertainty be encoded through regularizers, kernels, RKHS norms, and Gaussian-process priors?

\paragraph{Section map.}
Section 33.1 introduces regularization as explicit prior knowledge and derives ridge regression as both penalized least squares and MAP estimation. Section 33.2 defines kernels, feature maps, Gram matrices, and kernel ridge regression. Section 33.3 develops reproducing kernel Hilbert spaces and the representer theorem. Section 33.4 defines Gaussian processes as distributions over functions and derives the GP regression posterior.

\subsection{33.1 Regularization as prior knowledge}

\paragraph{Guiding question.}
How do we express a preference for simpler, smoother, smaller, or more structured functions inside an optimization problem?

\paragraph{Beginner-facing intuition.}
Empirical risk only asks how well a function fits the training data. Regularization adds a second question: among functions that fit reasonably well, which ones match our assumptions about the problem? Ridge regression prefers small coefficient norms. Lasso prefers sparse coefficients. Smoothness penalties prefer functions that do not wiggle too rapidly. Early stopping and dropout also act as regularizing mechanisms, though their exact mathematical forms differ by model.

Regularization is not an afterthought. It is one way inductive bias becomes a formula.

\paragraph{Formal definitions.}
Let \(\mathcal F=\{f_\theta:\theta\in\Theta\}\) be a model class and let
\[
\widehat R_n(\theta)=\frac1n\sum_{i=1}^n \ell(f_\theta(x_i),y_i)
\]
be empirical risk. A regularized objective has the form
\[
\boxed{
J_\lambda(\theta)=\widehat R_n(\theta)+\lambda\Omega(\theta),
}
\]
where \(\Omega(\theta)\geq0\) is a complexity penalty and \(\lambda\geq0\) controls its strength.

Common examples include
\[
\Omega(\theta)=\|\theta\|_2^2 \quad\text{(ridge or weight decay)}
\]
and
\[
\Omega(\theta)=\|\theta\|_1=\sum_j|\theta_j| \quad\text{(lasso sparsity)}.
\]
For functions, one may penalize a derivative norm such as \(\int |f'(x)|^2\,dx\), a roughness measure, or an RKHS norm introduced later in this chapter.

\paragraph{Core derivation: ridge regression and MAP.}
Consider linear regression with Gaussian noise:
\[
y=Xw+\varepsilon,\qquad \varepsilon\sim\mathcal N(0,\sigma^2 I_n).
\]
The negative log-likelihood, up to constants independent of \(w\), is
\[
\frac{1}{2\sigma^2}\|Xw-y\|_2^2.
\]
Place a zero-mean Gaussian prior on weights:
\[
w\sim\mathcal N(0,\tau^2 I_d).
\]
The negative log-prior, up to constants, is
\[
\frac{1}{2\tau^2}\|w\|_2^2.
\]
The negative log-posterior is therefore
\[
\frac{1}{2\sigma^2}\|Xw-y\|_2^2
+\frac{1}{2\tau^2}\|w\|_2^2
+\text{constant}.
\]
Multiplying by \(2\sigma^2\), the MAP estimator solves
\[
\boxed{
\widehat w_{\mathrm{ridge}}
\in\operatorname*{arg\,min}_w
\|Xw-y\|_2^2+\lambda\|w\|_2^2,
\qquad
\lambda=\frac{\sigma^2}{\tau^2}.
}
\]
Differentiating gives
\[
2X^\top(Xw-y)+2\lambda w=0,
\]
so
\[
\boxed{
(X^\top X+\lambda I)\widehat w=X^\top y.
}
\]
When \(\lambda>0\), the matrix \(X^\top X+\lambda I\) is invertible as long as no direction is completely unpenalized, even if \(X^\top X\) is singular.

\paragraph{Concrete example.}
Suppose \(X^\top X=9\), \(X^\top y=12\), and the one-dimensional ridge penalty is \(\lambda w^2\). Ordinary least squares gives
\[
\widehat w_{\mathrm{OLS}}=\frac{12}{9}=1.333.
\]
With \(\lambda=3\),
\[
\widehat w_{\mathrm{ridge}}=\frac{12}{9+3}=1.
\]
The coefficient is shrunk toward zero. This increases bias if the unregularized estimate was accurate, but it can reduce variance when the data are noisy or the design is ill-conditioned.

\paragraph{Computational neuroscience example.}
When fitting a receptive-field filter with many time lags and stimulus pixels, the number of coefficients can be large relative to the number of trials. A ridge penalty discourages enormous filter weights that fit trial noise. A smoothness penalty can encode the assumption that neighboring time lags or neighboring visual locations should have similar weights. A sparsity penalty can encode the assumption that only a few features drive the response. These are modeling assumptions, not biological facts by themselves.

\paragraph{AI and machine-learning example.}
Weight decay in neural networks is an \(L_2\)-style penalty on parameters. Data augmentation regularizes by making the learned function insensitive to transformations such as small translations or color changes. Dropout injects noise during training and can be interpreted as discouraging fragile co-adaptations. Different mechanisms express different prior beliefs about what should count as a stable predictor.

\paragraph{Common misconceptions and failure modes.}
First, regularization does not merely "make weights smaller"; it changes the optimization problem and the learned function. Second, a larger \(\lambda\) is not always better. Too much regularization underfits. Third, MAP is not the full Bayesian posterior; it is one point estimate. Fourth, a penalty only expresses useful prior knowledge if it matches the structure of the task.

\paragraph{Exercises.}
\begin{enumerate}[label=\textbf{33.1.\arabic*.}, leftmargin=*]
\item \textbf{Mechanical.} Write the ridge objective for \(X\in\mathbb R^{n\times d}\), \(y\in\mathbb R^n\), and \(w\in\mathbb R^d\). State the normal equations.
\item \textbf{Proof or derivation.} Starting from a Gaussian likelihood and Gaussian prior, derive the ridge objective and identify \(\lambda=\sigma^2/\tau^2\).
\item \textbf{Numerical/computational.} In one dimension, compute \(\widehat w_{\mathrm{ridge}}\) when \(X^\top X=20\), \(X^\top y=30\), and \(\lambda=10\). Compare with OLS.
\item \textbf{Interpretation.} Why might a smoothness penalty be more appropriate than an ordinary weight penalty for a visual receptive field?
\item \textbf{Misconception/debugging.} A student says, "Regularization always improves test error." Explain why too much regularization can hurt.
\end{enumerate}

\subsection{33.2 Kernel methods}

\paragraph{Guiding question.}
How can a method be linear in a high-dimensional feature space while producing nonlinear functions of the original input?

\paragraph{Beginner-facing intuition.}
A kernel is a similarity function that behaves like an inner product after mapping inputs into a feature space. Instead of explicitly constructing features \(\phi(x)\), a kernel method computes inner products \(k(x,x')\). This is the kernel trick: algorithms whose data dependence occurs through inner products can be run in an implicit feature space.

The result can be nonlinear in the original input even though it is linear in the feature space. Polynomial kernels create polynomial decision functions. Radial basis function kernels create smooth local similarity functions.

\paragraph{Formal definitions.}
Let \(\mathcal X\) be an input space. A feature map is a function
\[
\phi:\mathcal X\to\mathcal H,
\]
where \(\mathcal H\) is an inner-product space. The associated kernel is
\[
\boxed{
k(x,x')=\langle \phi(x),\phi(x')\rangle_{\mathcal H}.
}
\]
Given training inputs \(x_1,\ldots,x_n\), the Gram matrix is
\[
K\in\mathbb R^{n\times n},
\qquad
K_{ij}=k(x_i,x_j).
\]
A valid kernel must produce positive semidefinite Gram matrices:
\[
\boxed{
\sum_{i=1}^n\sum_{j=1}^n c_i c_j k(x_i,x_j)\geq0
\quad\text{for all }c\in\mathbb R^n.
}
\]
Equivalently, \(c^\top Kc\geq0\). This condition is automatic if \(k\) comes from an inner product because
\[
c^\top Kc
=\sum_{i,j}c_i c_j\langle\phi(x_i),\phi(x_j)\rangle
=\left\|\sum_i c_i\phi(x_i)\right\|_{\mathcal H}^2\geq0.
\]

\paragraph{Core mechanism: kernel ridge regression.}
Suppose a linear model is fitted in feature space:
\[
f_w(x)=\langle w,\phi(x)\rangle_{\mathcal H}.
\]
Kernel ridge regression solves
\[
\min_w \sum_{i=1}^n \bigl(\langle w,\phi(x_i)\rangle-y_i\bigr)^2+\lambda\|w\|_{\mathcal H}^2.
\]
The solution lies in the span of training features:
\[
w=\sum_{i=1}^n \alpha_i\phi(x_i).
\]
Then predictions on training inputs are
\[
f_w(x_j)=\sum_{i=1}^n\alpha_i k(x_j,x_i),
\]
or in vector form \(K\alpha\). The objective becomes
\[
\|K\alpha-y\|_2^2+\lambda\alpha^\top K\alpha.
\]
Differentiating and using the positive semidefinite structure gives the standard solution
\[
\boxed{
\alpha=(K+\lambda I)^{-1}y
}
\]
when the inverse exists. A new input \(x_*\) is predicted by
\[
\boxed{
\widehat f(x_*)=k_*^\top\alpha,
\qquad
k_*=(k(x_1,x_*),\ldots,k(x_n,x_*))^\top.
}
\]

\paragraph{Concrete example.}
Use the polynomial kernel
\[
k(x,x')=(1+xx')^2
\]
on scalar inputs \(x_1=0\), \(x_2=1\), \(x_3=2\). The Gram matrix is
\[
K=
\begin{bmatrix}
1 & 1 & 1\\
1 & 4 & 9\\
1 & 9 & 25
\end{bmatrix}.
\]
This kernel corresponds to an inner product in features proportional to \((1,\sqrt2 x,x^2)\), so a linear function in feature space is a quadratic function of \(x\). Kernel methods let the algorithm use that quadratic feature space through \(K\).

\paragraph{Computational neuroscience example.}
Suppose we want to predict firing rate from stimulus orientation. A radial basis function kernel
\[
k(x,x')=\exp\left(-\frac{(x-x')^2}{2\ell^2}\right)
\]
encodes the assumption that nearby orientations have similar responses. The length scale \(\ell\) controls smoothness: small \(\ell\) allows rapid changes, while large \(\ell\) enforces broad tuning. For circular variables such as orientation, a periodic kernel may be more appropriate than a Euclidean one.

\paragraph{AI and machine-learning example.}
Kernel classifiers and regressors can create nonlinear boundaries without manually engineering all nonlinear features. In modern large-scale AI, explicit kernel methods are sometimes too expensive for huge datasets, but kernel ideas remain central in similarity search, attention mechanisms, Gaussian processes, neural tangent kernels, and representation learning.

\paragraph{Common misconceptions and failure modes.}
First, any similarity function is not automatically a valid kernel; it must produce positive semidefinite Gram matrices. Second, the kernel defines the notion of similarity, so bad scaling or irrelevant features can make a good kernel perform poorly. Third, kernel methods still regularize; a flexible kernel with tiny regularization can overfit. Fourth, the kernel trick saves feature construction, not all computation: storing and inverting \(K\) can be expensive for large \(n\).

\paragraph{Exercises.}
\begin{enumerate}[label=\textbf{33.2.\arabic*.}, leftmargin=*]
\item \textbf{Mechanical.} For scalar inputs \(0,1\) and kernel \(k(x,x')=1+xx'\), write the \(2\times2\) Gram matrix.
\item \textbf{Proof or derivation.} Prove that if \(k(x,x')=\langle\phi(x),\phi(x')\rangle\), then every Gram matrix built from \(k\) is positive semidefinite.
\item \textbf{Numerical/computational.} With \(K=\begin{bmatrix}1&0\\0&1\end{bmatrix}\), \(y=(2,4)^\top\), and \(\lambda=1\), compute \(\alpha=(K+\lambda I)^{-1}y\).
\item \textbf{Interpretation.} In an RBF kernel for tuning curves, what does increasing the length scale \(\ell\) do to the fitted function?
\item \textbf{Misconception/debugging.} A student proposes \(k(x,x')=-|x-x'|\) because it is larger for nearby points. Explain why similarity intuition alone does not prove it is a valid kernel.
\end{enumerate}

\subsection{33.3 Reproducing kernel Hilbert spaces}

\paragraph{Guiding question.}
What kind of function space makes kernel evaluation, smoothness control, and finite-dimensional optimization fit together?

\paragraph{Beginner-facing intuition.}
A reproducing kernel Hilbert space, or RKHS, is a Hilbert space whose elements are functions and whose inner product is tied to point evaluation. The kernel \(k(x,\cdot)\) acts like a feature vector representing the point \(x\). The reproducing property says that evaluating a function at \(x\) is the same as taking an inner product with that representer.

This turns function fitting into geometry. A small RKHS norm means the function is simple according to the geometry defined by the kernel. Different kernels define different notions of simplicity.

\paragraph{Formal definitions.}
Let \(\mathcal X\) be an input set. A Hilbert space \(\mathcal H\) of functions \(f:\mathcal X\to\mathbb R\) is a reproducing kernel Hilbert space if there exists a function
\[
k:\mathcal X\times\mathcal X\to\mathbb R
\]
such that:
\[
k(x,\cdot)\in\mathcal H\quad\text{for every }x\in\mathcal X,
\]
and
\[
\boxed{
f(x)=\langle f,k(x,\cdot)\rangle_{\mathcal H}
\quad\text{for every }f\in\mathcal H.
}
\]
The second equation is the reproducing property. It says that point evaluation is a continuous linear functional on \(\mathcal H\).

The RKHS norm \(\|f\|_{\mathcal H}\) measures complexity according to the kernel. For many smooth kernels, functions with small RKHS norm are smooth in a kernel-specific sense. The exact smoothness meaning depends on the kernel and input domain.

\paragraph{Core theorem: representer theorem.}
Consider an optimization problem over an RKHS:
\[
\min_{f\in\mathcal H}
\sum_{i=1}^n L(f(x_i),y_i)
+\lambda\|f\|_{\mathcal H}^2,
\qquad \lambda>0.
\]
Assume \(L\) depends on \(f\) only through its values at the training inputs. Then every minimizer has the form
\[
\boxed{
f(\cdot)=\sum_{i=1}^n \alpha_i k(x_i,\cdot).
}
\]

\emph{Proof mechanism.} Let
\[
\mathcal S=\operatorname{span}\{k(x_1,\cdot),\ldots,k(x_n,\cdot)\}.
\]
Every \(f\in\mathcal H\) can be decomposed as
\[
f=f_{\parallel}+f_{\perp},
\]
where \(f_{\parallel}\in\mathcal S\) and \(f_{\perp}\perp\mathcal S\). For each training input \(x_i\),
\[
f(x_i)=\langle f,k(x_i,\cdot)\rangle_{\mathcal H}
=\langle f_{\parallel},k(x_i,\cdot)\rangle_{\mathcal H}
+\langle f_{\perp},k(x_i,\cdot)\rangle_{\mathcal H}.
\]
The second term is zero because \(k(x_i,\cdot)\in\mathcal S\) and \(f_{\perp}\perp\mathcal S\). Therefore \(f\) and \(f_{\parallel}\) have identical training predictions. The data-fit term is unchanged if we drop \(f_{\perp}\). But
\[
\|f\|_{\mathcal H}^2
=\|f_{\parallel}\|_{\mathcal H}^2+\|f_{\perp}\|_{\mathcal H}^2
\geq \|f_{\parallel}\|_{\mathcal H}^2.
\]
Thus any nonzero perpendicular component can only increase the regularizer without improving the fit. A minimizer can be chosen entirely in \(\mathcal S\), giving the finite expansion.

\paragraph{Concrete example.}
Let \(\mathcal X=\{a,b\}\) and suppose the kernel matrix is
\[
K=
\begin{bmatrix}
1 & 0.5\\
0.5 & 1
\end{bmatrix}.
\]
Any representer-theorem solution on training points \(a,b\) has the form
\[
f=\alpha_1 k(a,\cdot)+\alpha_2 k(b,\cdot).
\]
Its values on the training points are
\[
\begin{bmatrix}
f(a)\\
f(b)
\end{bmatrix}
=
\begin{bmatrix}
1 & 0.5\\
0.5 & 1
\end{bmatrix}
\begin{bmatrix}
\alpha_1\\
\alpha_2
\end{bmatrix}.
\]
Function fitting in the RKHS has become finite-dimensional coefficient fitting through the Gram matrix.

\paragraph{Computational neuroscience example.}
An RKHS model for a tuning curve can penalize functions that are rough under a chosen orientation kernel. The representer theorem says that the fitted curve is a weighted sum of kernel bumps centered at observed stimulus values. This is interpretable: the prediction at a new stimulus is assembled from similarities to measured stimuli. The choice of kernel encodes what "nearby" means for stimuli.

\paragraph{AI and machine-learning example.}
Kernel SVMs, kernel ridge regression, and many nonparametric regression methods rely on the representer theorem. The theorem explains why an infinite-dimensional function optimization can be solved using \(n\) coefficients, one per training example. It also clarifies a limitation: predictions often require comparing a new input to many training inputs unless approximation methods are used.

\paragraph{Common misconceptions and failure modes.}
First, an RKHS is not every possible function space; it is a Hilbert space where point evaluation has a kernel representer. Second, the RKHS norm is not always the same as an ordinary \(L_2\) norm over input space. Third, the representer theorem does not say all good functions have finite kernel expansions; it says regularized empirical objectives have minimizers of that form. Fourth, smooth-looking plots can hide poor kernel choices or bad length scales.

\paragraph{Exercises.}
\begin{enumerate}[label=\textbf{33.3.\arabic*.}, leftmargin=*]
\item \textbf{Mechanical.} State the reproducing property and identify the object \(k(x,\cdot)\).
\item \textbf{Proof or derivation.} Reproduce the orthogonal-decomposition proof of the representer theorem in your own words.
\item \textbf{Numerical/computational.} For \(K=\begin{bmatrix}2&1\\1&2\end{bmatrix}\) and \(\alpha=(1,-1)^\top\), compute \(f(x_1)\) and \(f(x_2)\).
\item \textbf{Interpretation.} Why does the representer theorem make kernel methods computationally finite even when the feature space is infinite-dimensional?
\item \textbf{Misconception/debugging.} A student says, "The RKHS norm is just the average squared value of the function." Explain why the norm is defined by the kernel-space inner product, not by ordinary averaging unless those structures coincide.
\end{enumerate}

\subsection{33.4 Gaussian processes}

\paragraph{Guiding question.}
How can we put a probability distribution on an unknown function and update that distribution after observing data?

\paragraph{Beginner-facing intuition.}
A Gaussian process, or GP, is a distribution over functions. Before seeing data, it says which functions are plausible through a mean function and a covariance kernel. After seeing noisy observations, Gaussian conditioning turns the prior over functions into a posterior over functions. The posterior mean is a prediction, and the posterior variance is uncertainty.

The key is finite-dimensional consistency. We never have to sample an infinite function all at once. For any finite set of input points, the corresponding function values form a multivariate Gaussian vector.

\paragraph{Formal definitions.}
A random function \(f:\mathcal X\to\mathbb R\) is a Gaussian process with mean function \(m\) and covariance kernel \(k\), written
\[
\boxed{
f\sim\mathcal{GP}(m,k),
}
\]
if for every finite collection \(x_1,\ldots,x_n\in\mathcal X\),
\[
\begin{bmatrix}
f(x_1)\\
\vdots\\
f(x_n)
\end{bmatrix}
\sim
\mathcal N
\left(
\begin{bmatrix}
m(x_1)\\
\vdots\\
m(x_n)
\end{bmatrix},
\begin{bmatrix}
k(x_1,x_1)&\cdots&k(x_1,x_n)\\
\vdots&\ddots&\vdots\\
k(x_n,x_1)&\cdots&k(x_n,x_n)
\end{bmatrix}
\right).
\]
The kernel must be positive semidefinite so that every finite covariance matrix is valid.

For regression, observations are usually modeled as
\[
y_i=f(x_i)+\varepsilon_i,
\qquad
\varepsilon_i\sim\mathcal N(0,\sigma_n^2),
\]
independently. Let
\[
y=(y_1,\ldots,y_n)^\top\in\mathbb R^n,
\]
\[
K\in\mathbb R^{n\times n},
\qquad
K_{ij}=k(x_i,x_j),
\qquad
I_n\in\mathbb R^{n\times n},
\]
and for one scalar test value \(f_*=f(x_*)\in\mathbb R\), let
\[
k_*=(k(x_1,x_*),\ldots,k(x_n,x_*))^\top\in\mathbb R^n
\]
be viewed as an \(n\times1\) column vector, and let
\[
k_{**}=k(x_*,x_*)\in\mathbb R.
\]

\paragraph{Core derivation: GP regression posterior.}
For simplicity, take \(m(x)=0\). The training observations satisfy
\[
y\sim\mathcal N(0,K+\sigma_n^2 I_n),
\]
where \(K+\sigma_n^2 I_n\in\mathbb R^{n\times n}\). The noise term adds variance to observed training outputs; it does not change the latent test variance \(k_{**}\).
The joint distribution of \(y\) and the latent test value \(f_*=f(x_*)\) is Gaussian:
\[
\begin{bmatrix}
y\\
f_*
\end{bmatrix}
\sim
\mathcal N
\left(
0,
\begin{bmatrix}
K+\sigma_n^2 I_n & k_*\\
k_*^\top & k_{**}
\end{bmatrix}
\right).
\]
The joint vector has shape \((n+1)\times1\). Its covariance has block dimensions
\[
\begin{bmatrix}
n\times n & n\times1\\
1\times n & 1\times1
\end{bmatrix},
\]
so the product \(k_*^\top(K+\sigma_n^2 I_n)^{-1}y\) is \(1\times1\), a scalar posterior mean for \(f_*\).
Gaussian conditioning gives
\[
\boxed{
f_*\mid y
\sim
\mathcal N
\left(
k_*^\top(K+\sigma_n^2 I_n)^{-1}y,\,
k_{**}-k_*^\top(K+\sigma_n^2 I_n)^{-1}k_*
\right).
}
\]
With a nonzero mean function, replace \(y\) by \(y-m_X\) and add \(m(x_*)\) to the posterior mean:
\[
\mu_*(x_*)=m(x_*)+k_*^\top(K+\sigma_n^2 I_n)^{-1}(y-m_X).
\]
For multiple test points \(x_*^{(1)},\ldots,x_*^{(m)}\), collect latent values in \(f_*\in\mathbb R^m\). Let \(K_{X*}\in\mathbb R^{n\times m}\) have entries \((K_{X*})_{ij}=k(x_i,x_*^{(j)})\), and let \(K_{**}\in\mathbb R^{m\times m}\) be the test-test kernel matrix. Then
\[
\mathbb E[f_*\mid y]=K_{X*}^\top(K+\sigma_n^2 I_n)^{-1}y\in\mathbb R^m,
\]
and
\[
\operatorname{Cov}(f_*\mid y)=K_{**}-K_{X*}^\top(K+\sigma_n^2 I_n)^{-1}K_{X*}\in\mathbb R^{m\times m}.
\]

This is a beginner-usable formula: build the noisy training covariance \(K+\sigma_n^2 I_n\), solve a linear system, compute a posterior mean, and compute a posterior covariance.

\paragraph{Concrete numerical example.}
Use a zero-mean GP with RBF kernel
\[
k(x,x')=\exp\left(-\frac{(x-x')^2}{2}\right)
\]
and noise variance \(\sigma_n^2=0.25\). Suppose one observation is \(x_1=0\), \(y_1=1\). At \(x_*=0\),
\[
K+\sigma_n^2 I_n=[1.25],\qquad k_*=[1].
\]
The posterior mean is
\[
\mu_*=1\cdot(1/1.25)\cdot1=0.8,
\]
and the posterior variance is
\[
1-1(1/1.25)1=0.2.
\]
At \(x_*=1\),
\[
k_*=[e^{-1/2}]\approx[0.607].
\]
The posterior mean is
\[
0.607(1/1.25)\approx0.486,
\]
and the posterior variance is
\[
1-\frac{0.607^2}{1.25}\approx0.706.
\]
The prediction is closer to the observed value near \(x=0\), and uncertainty increases farther away.

\paragraph{Computational neuroscience example.}
A GP can model an unknown tuning curve \(f(s)\) as a smooth random function of stimulus \(s\). Observed spike-rate estimates update the posterior. The posterior mean gives a fitted tuning curve, and the posterior variance identifies stimulus regions where data are sparse or noisy. GPs can also model latent trajectories in neural state space, though scaling and non-Gaussian spike observations often require approximations.

\paragraph{AI and machine-learning example.}
Gaussian processes are widely used in Bayesian optimization. The unknown function may be validation accuracy as a function of hyperparameters. The GP posterior mean estimates performance, while posterior uncertainty guides exploration. GP regression also clarifies the Bayesian interpretation of kernel ridge regression: for a zero-mean GP with Gaussian noise, the posterior mean has the same algebraic form as kernel ridge prediction with regularization tied to noise variance.

\paragraph{Practical caveats and misconceptions.}
First, a GP is not a distribution over parameter vectors; it is a distribution over function values with consistent Gaussian finite marginals. Second, the kernel is a modeling assumption. It determines smoothness, periodicity, amplitude, and extrapolation behavior. Third, exact GP regression requires solving systems involving an \(n\times n\) matrix, often costing \(O(n^3)\) time for direct factorization. Fourth, posterior variance is uncertainty under the model, not guaranteed real-world error. Misspecified kernels or shifted data can make it misleading.

\paragraph{Exercises.}
\begin{enumerate}[label=\textbf{33.4.\arabic*.}, leftmargin=*]
\item \textbf{Mechanical.} Define \(f\sim\mathcal{GP}(m,k)\) in terms of finite-dimensional Gaussian distributions.
\item \textbf{Proof or derivation.} Starting from the block joint Gaussian for \((y,f_*)\), write the conditional mean and variance formulas for \(f_*\mid y\).
\item \textbf{Numerical/computational.} Repeat the one-observation GP calculation above for noise variance \(\sigma_n^2=1\) at \(x_*=0\). What happens to the posterior mean and variance?
\item \textbf{Interpretation.} In GP tuning-curve fitting, what does a large posterior variance at a stimulus value usually indicate?
\item \textbf{Misconception/debugging.} A student says, "The GP posterior mean is the true function, and the variance is the true error bar." Correct this statement by distinguishing model-based posterior uncertainty from ground truth.
\end{enumerate}

\paragraph{Closing bridge.}
Regularization, kernels, RKHSs, and Gaussian processes show two complementary views of prior knowledge: as penalties in optimization and as priors in probability. Chapter 34 extends the probabilistic view from unknown functions to hidden causes, graphical factorization, latent variables, EM, and approximate posterior inference.

\clearpage
